\subsection{Perturbation}

In this number, E is a complex Hilbert space.

If \(u\) is a self-adjoint partial operator on E and \(v \in \mathcal{L}(\mathrm{E})\) is a Hermitian endomorphism, then \(u + v\) is self-adjoint (cf. IV, p. 236). This is not the case in general when \(v\) is a symmetric partial operator (exercise 9 of IV, p. 347). We shall nevertheless obtain positive results of this type in this number.

DEFINITION 11. --- Let u be a partial operator on E. A partial operator v on E is said to be bounded relative to u if dom(u) ⊂ dom(v) and if v defines, by restriction to the subspace, a continuous linear map from \(E_{u}\) into E.

Let u be a partial operator on E. Let v be a partial operator bounded relative to u. There exists a real number m such that

\[
\| v (x) \| \leqslant m (\| x \| + \| u (x) \|)
\]

for all \(x \in \text{dom}(u)\) . The relative norm of v with respect to u, denoted by \(\|v\|_{u}\), is the lower bound of the real numbers \(a \geqslant 0\) such that there exists a real number b such that

\[
\| v (x) \| \leqslant a \| u (x) \| + b \| x \|
\]

for all \(x \in \text{dom}(u)\) . The relative norm of v is therefore less than the norm of the restriction of v to the space \(E_{u}\) .

Note. --- Let u be a partial operator on E. Every endomorphism \(v \in \mathcal{L}(\mathrm{E})\) is bounded relative to u and its relative norm is zero since \(\|v(x)\| \leqslant \|v\|\|x\|\) for all \(x \in \operatorname{dom}(u)\) , which allows us to take a = 0 in the inequality above.

THEOREM 5 (Kato–Rellich). — Let u be a self-adjoint partial operator on E and v a symmetric partial operator bounded relative to u. If the relative norm \(\|v\|_{u}\) is \(< 1\), then the partial operator \(u + v\) whose domain is dom(u) is self-adjoint.

Since \(\| v \|_u < 1\) there exist, by definition, positive real numbers \(a\) and \(b\) such that \(a < 1\) and \(\| v(x) \| \leqslant a \| u(x) \| + b \| x \|\) for all \(x \in \text{dom}(u)\) .

Let \(t \in R^{*}\) . Set \(w_{t} = v \circ \mathrm{R}(u, it)\) . We have \(\operatorname{dom}(w_{t}) = E\) since the image of \(\mathrm{R}(u, it)\) is contained in the domain of u, which is contained in the domain of v by hypothesis. Let \(x \in \operatorname{dom}(u)\) . Since u is self-adjoint, we have

\[
\begin{array}{c} \| (i t - u) x \| ^ {2} = \| i t x \| ^ {2} + \| u (x) \| ^ {2} - i t \langle x | u (x) \rangle + i t \langle u (x) | x \rangle \\ = | t | ^ {2} \| x \| ^ {2} + \| u (x) \| ^ {2}, \end{array}
\]

whence the inequalities \(\|u(x)\|\leqslant\|(it-u)x\|\) and \(\|x\|\leqslant|t|^{-1}\|(it-u)x\|\) . But then it follows that

\[
\| v (x) \| \leqslant (a + b | t | ^ {- 1}) \| (i t - u) x \|.
\]

In particular, set \(x = \mathrm{R}(u, it)y\) for \(y \in \mathrm{E}\) ; we obtain

\[
\| w _ {t} (y) \| \leqslant (a + b | t | ^ {- 1}) \| y \|.
\]

We deduce that \(w_{t}\in\mathcal{L}(\mathrm{E})\) and that \(\|w_{t}\|\leqslant a+|t|^{-1}b\) .

Since \(a < 1\) , there exists \(t \in \mathbf{R}_{+}^{*}\) such that \(a + b|t|^{-1} < 1\) . We then have \(\|w_{t}\| < 1\) and \(\|w_{-t}\| < 1\) ; the endomorphisms \(1_{\mathrm{E}} - w_{t}\) and \(1_{\mathrm{E}} - w_{-t}\) of E are therefore invertible (prop. 2 of I, p. 22). Since we have the formula \((1_{\mathrm{E}} - w_{t}) \circ (it - u) = it - (u + v)\) , this implies that \(u + v - it\) is surjective; likewise, the partial operator \(u + v + it\) is surjective. It follows that \(u + v\) is self-adjoint (prop. 11 of IV, p. 240).

